\section{Introduction}

% ¶1 具体问题切入（ATC风格：快速进入实际场景）
Mixture-of-Experts (MoE) models scale language model capacity by routing each token through only a few selected experts, keeping per-token computation manageable while expanding the total parameter count to hundreds of billions~\cite{lepikhin2021gshard,fedus2022switch,rajbhandari2022deepspeed}.
This efficiency makes MoE the architecture of choice for many open-weight frontier models, including Qwen3, Mixtral, and DeepSeek-V3.
For practitioners serving these models on commodity GPUs, however, the large number of experts creates a tension:
the aggregate expert weight far exceeds GPU memory, forcing expert weights to be offloaded to host DRAM.
Speculative decoding (SD) is the standard technique for accelerating autoregressive LLM decoding, routinely delivering $2$--$3\times$ speedup on dense models~\cite{leviathan2023fast,chen2023accelerating,li2024eagle,li2025eagle3}.
A natural question is whether SD can similarly accelerate offloaded MoE serving.
We find that, without careful system support, \emph{it cannot}.

% ¶2 Motivating Example + Figure 1
Consider deploying Qwen3-30B-A3B on a 48\,GB GPU.
Its 128 experts per layer require about 57.6\,GB in BF16---exceeding the GPU budget---so the standard remedy is CPU offloading via PCIe.
As Figure~\ref{fig:motivation} shows, naively enabling SD under this setting yields a mere $1.17\times$ speedup, far below the $2$--$3\times$ expected for dense models.
The culprit is not a weak draft model.
Multi-token verification forces each layer to activate a wider set of experts: $\wsd \approx 19.3$ on average (up to 26.9), compared with $\war = 8$ for single-token decoding, amplifying PCIe transfer volume by $2.4\times$.
SD fails here not because it lacks algorithmic advantage, but because it \emph{amplifies the memory bottleneck}.

\begin{figure}[t]
\centering
\includegraphics[width=\columnwidth]{figures/fig_motivation.pdf}
\caption{Profiling Qwen3-30B-A3B on one 48\,GB GPU. (a)~SD expands the per-layer expert working set from $\war=8$ to $\wsd\approx 19.3$ (up to 26.9), amplifying PCIe transfer by $2.4\times$. (b)~End-to-end throughput: SD alone yields only $1.17\times$, but \sysname reaches $4.80\times$ by eliminating cache-path bottlenecks.}
\Description{Panel a shows a bar chart of 48 MoE layers comparing AR working set of 8 with SD working set ranging from 14 to 27. Panel b shows a 4-way grouped bar chart at 44 GB where SD plus BriskMoE reaches 9.98 tokens per second.}
\label{fig:motivation}
\end{figure}

% ¶3 缓存——必要但远不充分（核心洞见段）
A well-known technique to mitigate offloading overhead is to cache frequently used experts in GPU memory, as done by MoE-Infinity~\cite{xue2024moeinfinity} and pre-gated MoE~\cite{hwang2023pregated}.
We derive a minimum GPU memory budget $M^*$ required for the cache to accommodate SD's expanded working set, and confirm that allocating at least $M^*$ is a necessary precondition (Section~\ref{sec:model}).
However, we find that \emph{caching is necessary but far from sufficient}: even when every expert access hits the GPU cache, the system achieves only 29\% of peak compute efficiency.
Our profiling reveals two previously unrecognized bottlenecks that are unique to the combination of SD and expert caching:
%
\begin{itemize}[leftmargin=1.2em,topsep=2pt,itemsep=1pt]
\item \textbf{Cache-hit path overhead.}
Even on a full cache hit, scratchpad memory copies and GPU--CPU synchronization on the cache-hit path consume $2.4\times$ the actual kernel compute time~(\textbf{Obs.\,1}).
This overhead exists regardless of cache hit rate; prior cached-MoE systems never exposed it because AR decoding processes only one token per step, making the absolute expert compute---and thus its overhead---negligible relative to other latency components such as attention.
SD's multi-token verification turns expert computation into the critical path, amplifying this previously hidden cost into the dominant bottleneck.

\item \textbf{SD-induced transient overflow.}
SD's per-step working set fluctuates well above the steady-state average, triggering bursts of cache evictions that cascade across $\sim$8 subsequent steps~(\textbf{Obs.\,2}).
Cache-based systems designed for AR decoding do not anticipate this transient pressure because single-token inference produces a stable, predictable access pattern.
\end{itemize}

% ¶4 BriskMoE 方案 + 实用性
To address these bottlenecks, we present \sysname, a lightweight system that makes speculative decoding an effective latency accelerator for offloaded MoE inference on a single GPU.
\sysname targets each bottleneck with a dedicated design:
(1)~\textbf{Unblock} eliminates hit-path overhead through direct in-pool expert execution, synchronization-free cache routing, cross-layer prefetching, and fused-MoE kernel tuning;
(2)~\textbf{Protect} prevents SD-induced cache overflow by exploiting the draft-generation interval to pre-stage the experts that verification will likely need.
\sysname is implemented as a non-intrusive plugin for vLLM~\cite{kwon2023vllm} with fewer than 800 lines of core logic and requires zero modification to the inference runtime, making it immediately deployable on existing serving stacks.

% ¶5 贡献列举（动词开头 + section引用）
In summary, this paper makes the following contributions.
\begin{itemize}[leftmargin=1.2em]
\item We present a systematic profiling study of speculative decoding under single-request offloaded MoE serving. We derive a minimum memory threshold $M^*$ that determines whether SD can be effective and identify two previously unrecognized bottlenecks---cache-hit path overhead and SD-induced transient overflow---that explain why expert caching alone is insufficient (Section~\ref{sec:model}).

\item We develop \sysname, a system with two complementary designs---\emph{Unblock} and \emph{Protect}---that systematically removes the identified bottlenecks. \sysname is implemented as a non-intrusive vLLM plugin, requiring zero changes to the host runtime (Sections~\ref{sec:design} and~\ref{sec:impl}).

\item We evaluate \sysname on two MoE architectures---Qwen3-30B-A3B (BF16, top-8) and GPT-OSS-120B (MXFP4, top-4)---across three workloads. \sysname achieves up to $4.80\times$ speedup on Qwen3-30B and $5.00\times$ on GPT-OSS-120B over the offload baseline, restoring SD to a $1.36\times$ multiplier over cached autoregressive decoding (Section~\ref{sec:eval}).
\end{itemize}

% ¶6 论文组织
The rest of the paper is organized as follows. Section~\ref{sec:background} introduces the serving scenario and notation. Section~\ref{sec:model} presents the profiling study and analytical model. Section~\ref{sec:design} describes the \sysname design. Section~\ref{sec:impl} details the implementation. Section~\ref{sec:eval} reports experimental results. Section~\ref{sec:related} discusses related work and Section~\ref{sec:conclusion} concludes.
